\subsubsection{Why Variational Algorithms Matter}\label{sec:why-variational-algorithms-matter}

After studying all the strengths and weaknesses, a natural question arises: \emph{\textbf{if VQAs are approximate, slow, and suffer from barren plateaus... why do we study them at all?}} The answer is because they are much more than just an algorithm. They are also a \textbf{research paradigm} for learning how to exploit today's quantum computers.

\highspace
Variational Quantum Algorithms (VQAs) are important because they:
\begin{enumerate}
    \item \hl{Studying the Computational Power of Quantum Devices.} Today's quantum computers are still experimental. We still do not fully know:
    \begin{itemize}
        \item Which problems they solve efficiently;
        \item Which circuit structures work best;
        \item How hardware behaves under realistic workloads.
    \end{itemize}
    VQAs provide an ideal testing framework. So, instead of asking ``\emph{can this quantum computer factor 2048-bit integers?}'', we ask ``\emph{how well can it optimize a small Hamiltonian?}'' or ``\emph{how much entanglement can it reliably generate?}''. VQAs therefore serve as \textbf{benchmarks} for quantum hardware.
    
    
    \item \hl{Developing Methods for Noisy Quantum Computers.} The problem with NISQ devices is that they have a limited number of qubits and significant noise sources such as decoherence and gate errors. If we waited for perfect quantum computers, research would almost stop. Instead, \hl{VQAs encourage researchers to develop techniques specifically for imperfect devices.} Many techniques that are now standard in quantum computing (e.g., error mitigation, circuit compilation, and noise-aware optimization) were first explored within the VQA framework.
    
    So VQAs are not only useful for solving problems, they are also a laboratory for developing new quantum computing techniques.
    
    
    \item \hl{They May Still Be Useful in the Future.} VQAs may still be useful in the future for problems for which we do not yet have a dedicated quantum algorithm. Probably the optimizer will need to be part of the quantum circuit to ensure scalability.
    
    \highspace
    Today, the optimizer is classical: quantum circuits are run on a quantum computer, and the results are fed into a classical optimizer to update the parameters. This works for today's relatively small problems.

    \highspace
    But imagine a future quantum computer with millions of logical qubits. The optimizer may need to adjust millions of parameters. Sending information back and forth between quantum and classical computers could become the bottleneck. The communication overhead and the classical optimization itself may no longer scale well.

    \highspace
    One possible future solution is to implement the \textbf{quantum circuit} and the \textbf{optimizer} on the \textbf{same quantum computer}. The idea is that part (or all) of the optimization could itself be performed quantum mechanically. However, \hl{even if this is possible, many optimization problems may still not have a dedicated quantum algorithm and designing a specialized algorithm for every problem is unrealistic.} So, \textbf{VQAs may still be the best option for many problems in the future}.
\end{enumerate}
In summary, Variational Quantum Algorithms (VQAs) are important because they allow researchers to study the computational capabilities of current quantum devices and develop techniques to cope with limited resources and noise. Although they were originally motivated by the NISQ era, they may remain useful in the future for optimization problems that lack dedicated quantum algorithms. As quantum hardware scales, the optimization process itself may eventually need to become quantum to ensure scalability.